% ! TeX program = lualatex
% Reading one-sided limits off a graph.  Used by the midterm review slides.
%
%   page 1 -- a jump at x = 2.  The inputs climb to 2 from the left and the
%             outputs climb to 3, but the right-hand branch starts somewhere
%             else, so the two-sided limit does not exist.
%   page 2 -- the same reading from both sides, agreeing on 3, so the
%             two-sided limit does exist.
%
% The branches are built from square roots because a square root approaches
% its endpoint slowly: inputs spaced out along the x-axis give outputs that
% stay visibly separated instead of piling onto the limit at once.
%
% Three colours, three roles: the inputs sweep along the x-axis in main, the
% dots they land on are attn, and the outputs they produce climb the y-axis
% in supp.
\documentclass[multi=tikzpicture,crop]{standalone}

\input{../typography.tex.preamble}
\input{../tikz.tex.preamble}
\usetikzlibrary{arrows.meta}

\pgfplotsset{left limit/.style={
    basic curve,
    enlargelimits=false,
    xmin=-0.8, xmax=4.4,
    ymin=-0.5, ymax=3.9,
    xtick distance={1}, ytick distance={1},
    % the pictures are scaled down on the slide, so the title is a size up
    title style={font=\footnotesize, yshift=4pt},
    width={210pt},
  }}

% The left-hand branch is shared by both pictures.
\newcommand{\leftbranch}{%
  \addplot[very thick, domain=0:2, samples=100] {3 - sqrt(2-x)};
}

% The tracing: inputs along the x-axis, the dots they land on, and the
% outputs they produce.  Only the first picture carries it.
\newcommand{\tracing}{%
  \addplot[attn, only marks, mark=*, mark size=1.7pt] coordinates
  {(0.5,1.7753) (1,2) (1.5,2.2929) (1.75,2.5) (1.9,2.6838)};
  \draw[main, very thick, -{Stealth}] (axis cs:0.4,0) -- (axis cs:1.95,0);
  \draw[supp, very thick, -{Stealth}] (axis cs:0,1.7) -- (axis cs:0,2.9);
}

\begin{document}

% ---------------------------------------------------------------------
% Page 1 -- the right-hand branch jumps, so the two-sided limit fails.
\begin{tikzpicture}
  \begin{axis}[left limit,
      title={\(\lim_{x \to 2^{-}} f(x) = 3\) but \(\lim_{x \to 2} f(x)\) DNE}]
    \leftbranch{}
    \addplot[very thick, domain=2:4] {1.5 + (x-2)/2};
    \draw[fill=white] (axis cs:2,3) circle (0.06);
    \draw[fill=white] (axis cs:2,1.5) circle (0.06);
    \tracing{}
  \end{axis}
\end{tikzpicture}

% ---------------------------------------------------------------------
% Page 2 -- the right-hand branch mirrors the left, so both sides agree on 3.
\begin{tikzpicture}
  \begin{axis}[left limit, title={\(\lim_{x \to 2} f(x) = 3\)}]
    \leftbranch{}
    \addplot[very thick, domain=2:4, samples=100] {3 - sqrt(x-2)};
    \draw[fill=white] (axis cs:2,3) circle (0.06);
  \end{axis}
\end{tikzpicture}
\end{document}
